\documentclass[12pt,a4paper]{article}

% ── Codificação e idioma ─────────────────────────────────────────────
\usepackage[utf8]{inputenc}
\usepackage[T1]{fontenc}
\usepackage[portuguese]{babel}
\usepackage{lmodern}
\usepackage{microtype}

% ── Layout ───────────────────────────────────────────────────────────
\usepackage[top=2.8cm,bottom=2.8cm,left=3.0cm,right=2.5cm,headheight=16pt]{geometry}

% ── Cores ────────────────────────────────────────────────────────────
\usepackage[dvipsnames,table]{xcolor}
\definecolor{navy}{HTML}{1A3A5C}
\definecolor{medblue}{HTML}{2E75B6}
\definecolor{ltblue}{HTML}{D6E4F0}
\definecolor{rowalt}{HTML}{EBF3FB}
\definecolor{legalbg}{HTML}{F0FDF4}
\definecolor{legalbdr}{HTML}{16A34A}
\definecolor{alertbg}{HTML}{FFF8E1}
\definecolor{alertbdr}{HTML}{E59C0A}
\definecolor{warnbg}{HTML}{FEF2F2}
\definecolor{warnbdr}{HTML}{DC2626}
\definecolor{infobg}{HTML}{EFF6FF}
\definecolor{infobdr}{HTML}{3B82F6}
\definecolor{corrbg}{HTML}{FFF0F6}
\definecolor{corrbdr}{HTML}{9333EA}

% ── Cabeçalho / rodapé ───────────────────────────────────────────────
\usepackage{fancyhdr,lastpage}
\pagestyle{fancy}\fancyhf{}
\fancyhead[L]{\footnotesize\color{navy}\textbf{ATP --- Agenciamento Fluvial}
  --- Processo 2026.00002 --- Borba/AM}
\fancyhead[R]{\footnotesize\color{navy}Lei n.\textsuperscript{o} 14.133/2021}
\fancyfoot[C]{\footnotesize Página \thepage\ de \pageref{LastPage}}
\renewcommand{\headrulewidth}{0.5pt}
\renewcommand{\headrule}{%
  \hbox to\headwidth{\color{medblue}\leaders\hrule height 0.5pt\hfill}}

% ── Caixas coloridas ─────────────────────────────────────────────────
\usepackage{tcolorbox}
\tcbuselibrary{skins,breakable}

\newtcolorbox{legalbox}[1][Base Legal]{enhanced,breakable,
  colback=legalbg,colframe=legalbdr,boxrule=0.8pt,arc=4pt,
  left=10pt,right=10pt,top=5pt,bottom=5pt,
  title=\textbf{#1},fonttitle=\small\bfseries\color{legalbdr},
  attach boxed title to top left={yshift=-2mm,xshift=6mm},
  boxed title style={colback=legalbg,colframe=legalbdr,arc=3pt,boxrule=0.6pt}}

\newtcolorbox{alertbox}[1][Atenção]{enhanced,breakable,
  colback=alertbg,colframe=alertbdr,boxrule=0.8pt,arc=4pt,
  left=10pt,right=10pt,top=5pt,bottom=5pt,
  title=\textbf{#1},fonttitle=\small\bfseries\color{navy},
  attach boxed title to top left={yshift=-2mm,xshift=6mm},
  boxed title style={colback=alertbdr!25,colframe=alertbdr,arc=3pt,boxrule=0.6pt}}

\newtcolorbox{warnbox}[1][Risco]{enhanced,breakable,
  colback=warnbg,colframe=warnbdr,boxrule=0.8pt,arc=4pt,
  left=10pt,right=10pt,top=5pt,bottom=5pt,
  title=\textbf{#1},fonttitle=\small\bfseries\color{warnbdr},
  attach boxed title to top left={yshift=-2mm,xshift=6mm},
  boxed title style={colback=warnbg,colframe=warnbdr,arc=3pt,boxrule=0.6pt}}

\newtcolorbox{infobox}[1][Informação]{enhanced,breakable,
  colback=infobg,colframe=infobdr,boxrule=0.8pt,arc=4pt,
  left=10pt,right=10pt,top=5pt,bottom=5pt,
  title=\textbf{#1},fonttitle=\small\bfseries\color{infobdr},
  attach boxed title to top left={yshift=-2mm,xshift=6mm},
  boxed title style={colback=infobg,colframe=infobdr,arc=3pt,boxrule=0.6pt}}

\newtcolorbox{corrbox}[1][Correção Aplicada]{enhanced,breakable,
  colback=corrbg,colframe=corrbdr,boxrule=0.8pt,arc=4pt,
  left=10pt,right=10pt,top=5pt,bottom=5pt,
  title=\textbf{#1},fonttitle=\small\bfseries\color{corrbdr},
  attach boxed title to top left={yshift=-2mm,xshift=6mm},
  boxed title style={colback=corrbg,colframe=corrbdr,arc=3pt,boxrule=0.6pt}}

% ── Títulos ──────────────────────────────────────────────────────────
\usepackage{titlesec}
\titleformat{\section}{\color{navy}\large\bfseries}{\thesection.}{0.6em}{}
  [\color{navy}\titlerule]
\titlespacing{\section}{0pt}{22pt}{8pt}
\titleformat{\subsection}{\color{medblue}\normalsize\bfseries}{\thesubsection}{0.6em}{}
\titlespacing{\subsection}{0pt}{14pt}{5pt}
\titleformat{\subsubsection}{\color{navy!80}\normalsize\bfseries\itshape}
  {\thesubsubsection}{0.5em}{}
\titlespacing{\subsubsection}{0pt}{10pt}{4pt}

% ── Tabelas ──────────────────────────────────────────────────────────
\usepackage{booktabs,array,tabularx,longtable,multirow,colortbl}
\renewcommand{\arraystretch}{1.40}
\newcolumntype{C}[1]{>{\centering\arraybackslash}p{#1}}
\newcolumntype{L}[1]{>{\raggedright\arraybackslash}p{#1}}
\newcolumntype{R}[1]{>{\raggedleft\arraybackslash}p{#1}}

% ── Utilitários ──────────────────────────────────────────────────────
\usepackage{float,enumitem,setspace,amsmath,caption,chngcntr,pdflscape}
\setlist[enumerate]{leftmargin=1.8em,itemsep=3pt,topsep=3pt,parsep=0pt}
\setlist[itemize]{leftmargin=1.8em,itemsep=2pt,topsep=3pt,parsep=0pt}
\captionsetup{font=small,labelfont=bf,skip=4pt}
\counterwithout{footnote}{section}
\usepackage[colorlinks=true,linkcolor=navy,urlcolor=medblue,
  pdfborder={0 0 0}]{hyperref}

% ── Macros ───────────────────────────────────────────────────────────
\newcommand{\art}[1]{art.\,#1}
\newcommand{\Rs}{R\$\,}

% ── Dados do processo ─────────────────────────────────────────────────
\newcommand{\procNum}{2026.00002}
\newcommand{\exercicio}{2026}
\newcommand{\municipio}{Prefeitura Municipal de Borba/AM}
\newcommand{\valorTotal}{2.800.000,00}

% ═════════════════════════════════════════════════════════════════════
\begin{document}
% ═════════════════════════════════════════════════════════════════════

% ── Capa ─────────────────────────────────────────────────────────────
\begin{titlepage}
  \pagecolor{navy}\color{white}\setlength{\parindent}{0pt}\centering
  \vspace*{0.8cm}
  {\small\bfseries\color{ltblue}PREFEITURA MUNICIPAL DE BORBA ---
    ESTADO DO AMAZONAS}\\[2pt]
  {\small\color{ltblue!70}SECRETARIA MUNICIPAL DE SAÚDE --- PROCESSO
    \procNum}\\[0.4cm]
  \textcolor{ltblue!50}{\rule{0.85\linewidth}{0.8pt}}\\[0.5cm]
  {\fontsize{11}{14}\selectfont\bfseries\color{ltblue!85}
    ANÁLISE TÉCNICA PRELIMINAR --- ATP}\\[0.3cm]
  {\fontsize{11}{14}\selectfont\color{ltblue!70}
    Fundamento para o Estudo Técnico Preliminar (ETP) ---
    IN SEGES/MGI n.\textsuperscript{o} 58/2022}\\[0.5cm]
  {\fontsize{20}{26}\selectfont\bfseries
    Contratação de Empresa\\[4pt]Especializada em
    Agenciamento\\[4pt]de Transporte Fluvial\\[4pt]de Passageiros}\\[0.3cm]
  {\fontsize{12}{16}\selectfont\color{ltblue!80}
    Corredor Borba--Manaus (Rio Madeira) --- \exercicio}\\[0.6cm]
  \textcolor{ltblue!50}{\rule{0.85\linewidth}{0.8pt}}\\[0.8cm]
  \begin{tabular}{@{}r@{\hspace{8pt}}p{0.56\linewidth}@{}}
    \color{ltblue!70}\small\bfseries Processo: &
      \small \procNum \\[4pt]
    \color{ltblue!70}\small\bfseries Objeto: &
      \small Serviço de agenciamento de passagens fluviais \\[4pt]
    \color{ltblue!70}\small\bfseries Natureza da despesa: &
      \small 33.90.39 --- Outros Serviços de Terceiros PJ \\[4pt]
    \color{ltblue!70}\small\bfseries CATSER principal: &
      \small 3719 --- Agenciamento de Passagens \\[4pt]
    \color{ltblue!70}\small\bfseries Valor estimado: &
      \small \Rs \valorTotal\ (dois milhões e oitocentos mil reais) \\[4pt]
    \color{ltblue!70}\small\bfseries Modalidade proposta: &
      \small Pregão Eletrônico --- SRP (Lei n.\textsuperscript{o}
        14.133/2021) \\[4pt]
    \color{ltblue!70}\small\bfseries Vigência projetada: &
      \small 12 meses prorrogáveis (\art{84}, \S\,1.\textsuperscript{o}) \\[4pt]
    \color{ltblue!70}\small\bfseries Elaborado em: &
      \small Borba/AM, \exercicio \\[4pt]
  \end{tabular}
  \vfill
  \textcolor{ltblue!40}{\rule{0.85\linewidth}{0.4pt}}\\[0.3cm]
  {\tiny\color{ltblue!50}
    Documento elaborado em conformidade com o \art{18} da Lei
    n.\textsuperscript{o} 14.133/2021 e com a IN SEGES/MGI n.\textsuperscript{o}
    58/2022, com fundamento no princípio da eficiência (\art{37}, CF/88) e no
    planejamento da contratação}
\end{titlepage}
\nopagecolor
\setlength{\parindent}{1.2cm}

\tableofcontents
\clearpage

% ═════════════════════════════════════════════════════════════════════
\section{Precisão Jurídica do Objeto: Agenciamento \textit{vs.}
  Aquisição de Passagens}
% ═════════════════════════════════════════════════════════════════════

\begin{legalbox}[Distinção Fundamental --- \art{6.\textsuperscript{o}},
  XXIII, ``a'', Lei n.\textsuperscript{o} 14.133/2021]
\small
O objeto desta contratação é um \textbf{serviço continuado}, nos termos do
\art{6.\textsuperscript{o}}, XVI, da Lei n.\textsuperscript{o}~14.133/2021.
Não se trata de compra de bem (passagem) nem de contratação direta junto às
armadoras, mas de \textbf{intermediação profissional especializada}.
\end{legalbox}

A precisão do objeto é o alicerce de toda a estrutura licitatória. A tentação
de classificar este processo como ``aquisição de passagens fluviais'' produz
três consequências deletérias: (i) erro na natureza da despesa orçamentária,
que deve ser \textbf{33.90.39} (Outros Serviços de Terceiros --- Pessoa Jurídica),
e não 33.90.30 (Material de Consumo); (ii) inadequação do CATSER, que é
\textbf{3719 --- Agenciamento de Passagens}, não um código de produto; e (iii)
estrutura contratual incorreta --- a passagem é o \textit{insumo} do serviço,
não o objeto do contrato.

\begin{table}[H]\centering\small
\caption{Distinção jurídica entre agenciamento e aquisição direta}
\rowcolors{2}{rowalt}{white}
\begin{tabular}{L{4.0cm} L{5.0cm} L{5.0cm}}
  \rowcolor{navy}
  \color{white}\textbf{Dimensão} &
  \color{white}\textbf{Agenciamento (objeto correto)} &
  \color{white}\textbf{Aquisição direta (objeto incorreto)} \\
  \toprule
  Natureza jurídica       & Prestação de serviço continuado
                            (\art{6.\textsuperscript{o}}, XVI) &
                            Compra de bem (\art{6.\textsuperscript{o}}, XIV) \\
  Natureza da despesa     & 33.90.39 & 33.90.30 \\
  CATSER/CATMAT           & CATSER 3719 &
                            Não há CATMAT para passagens fluviais \\
  Instrumento contratual  & Contrato de serviço + ARP &
                            Apenas ARP ou empenho direto \\
  Quem remunera o armador & A agência, com seus recursos;
                            a prefeitura paga a agência &
                            A prefeitura diretamente ao armador \\
  Vantagem para a gestão  & Fatura unificada, bilhetes eletrônicos,
                            controle por centro de custo &
                            Reembolsos pulverizados, risco de superfaturamento \\
  \bottomrule
\end{tabular}
\end{table}

\begin{corrbox}[Correção 1 --- DPVAT $\to$ DPEM]
\small
O estudo de origem menciona o \textbf{DPVAT} como seguro obrigatório das
embarcações. Este seguro foi \textbf{extinto pela Lei n.\textsuperscript{o}
14.071, de 13 de outubro de 2020}. Para embarcações, o seguro obrigatório
vigente é o \textbf{DPEM --- Dano Pessoal Causado por Embarcações ou por
sua Carga}, regulado pela \textbf{Resolução CNSP n.\textsuperscript{o}
373/2018}, que estabelece os valores mínimos de indenização por morte e
invalidez permanente de passageiro. O TR deverá exigir a certidão de apólice
DPEM vigente como requisito habilitatório das embarcações cadastradas pela
agência.
\end{corrbox}

\begin{corrbox}[Correção 2 --- Natureza da Contratação e Estrutura de Lotes]
\small
O estudo de origem estrutura três ``itens'' (passagem fluvial, agenciamento e
frete de carga) como se fossem itens paralelos de igual natureza. A estrutura
correta é:
\begin{itemize}[topsep=2pt,itemsep=1pt]
  \item \textbf{Lote único de serviço}: agenciamento de passagens fluviais
    (CATSER 3719) com os lances incidindo sobre a \textbf{taxa de agenciamento
    por bilhete emitido};
  \item O frete de volumes acompanhantes (\textit{bagagem de porão}) pode
    ser item separado dentro do mesmo lote (CATSER 17892), desde que o TR
    justifique o agrupamento por unidade funcional (\art{40},
    \S\,1.\textsuperscript{o}, II, Lei n.\textsuperscript{o}~14.133/2021);
  \item \textbf{O preço da passagem não é licitado}: é preço de mercado
    regulado pela ARSEP-AM e repassado ao Município sem margem; a competição
    ocorre exclusivamente sobre a taxa de agenciamento.
\end{itemize}
\end{corrbox}

% ═════════════════════════════════════════════════════════════════════
\section{Fundamentos Matemáticos e Estruturação Orçamentária}
% ═════════════════════════════════════════════════════════════════════

\subsection{Custo Total do Serviço e Capacidade Operacional de Transporte}

O orçamento de \Rs \valorTotal\ representa o \textbf{teto de empenho} da Ata
de Registro de Preços e deve ser decomposto no \textbf{Custo Total do Serviço
de Agenciamento} ($CTSA$):

\begin{equation}
  CTSA = \underbrace{N_{bilhetes} \times \bar{P}}_{\text{repasse ao armador}}
       + \underbrace{N_{bilhetes} \times T_{ag}}_{\text{remuneração da agência}}
       + \underbrace{N_{kg} \times P_{frete}}_{\text{frete de volumes}}
  \label{eq:ctsa}
\end{equation}

onde $N_{bilhetes}$ é o número de passagens emitidas no período, $\bar{P}$
a média ponderada das tarifas ARSEP-AM por modal, $T_{ag}$ a taxa de
agenciamento por bilhete e $N_{kg} \times P_{frete}$ o custo do serviço
de despacho de carga associado.

\medskip
A \textbf{Capacidade Operacional de Transporte} ($COT$), derivada diretamente
do orçamento disponível, é dada por:

\begin{equation}
  COT = \frac{O_{total} - C_{frete}}{\bar{P} + T_{ag}}
  \label{eq:cot}
\end{equation}

onde $C_{frete}$ é a reserva orçamentária para frete de volumes (estimada em
5\% do total, conforme histórico de contratos similares na SEFAZ-AM):

\begin{align}
  C_{frete}   &= 0{,}05 \times 2{.}800{.}000 = \Rs 140{.}000 \notag\\
  O_{passagens} &= 2{.}800{.}000 - 140{.}000 = \Rs 2{.}660{.}000 \notag
\end{align}

\begin{table}[H]\centering
\caption{Projeção estática de volume por modal de transporte (orçamento
  para passagens: \Rs 2.660.000)}
\small\rowcolors{2}{rowalt}{white}
\begin{tabular}{L{4.5cm} C{2.0cm} C{2.0cm} C{2.2cm} C{2.5cm}}
  \rowcolor{navy}
  \color{white}\textbf{Modalidade} &
  \color{white}\textbf{Tarifa\newline ARSEP (R\$)} &
  \color{white}\textbf{Taxa\newline Ag. (R\$)} &
  \color{white}\textbf{CUF (R\$)} &
  \color{white}\textbf{Cap. Máx.\newline (bilhetes)} \\
  \toprule
  Lancha rápida (A-Jato) & 335,00 & 0,01 & 335,01 & 7.940 \\
  Barco motor (Recreio)  & 185,00 & 0,01 & 185,01 & 14.378 \\
  Misto (60\% recreio / 40\% jato) &
    245,00 & 0,01 & 245,01 & 10.856 \\
  \rowcolor{ltblue}
  \multicolumn{4}{l}{\textbf{Cenário base adotado (60/40)}} &
    \textbf{10.856 bilhetes} \\
  \bottomrule
\end{tabular}
\end{table}

\noindent O mix 60/40 (recreio/jato) reflete o perfil histórico de uso
institucional: pacientes de TFD e servidores em deslocamento rotineiro
preferem o recreio pelo custo e conforto em viagens diurnas; urgências
médicas e deslocamentos de gestores utilizam o A-Jato.

\subsection{Análise de Inexequibilidade da Taxa de Agenciamento}

\begin{warnbox}[Risco Jurídico --- Taxa Simbólica]
\small
A taxa de \Rs 0,01 por bilhete, embora documentada no PE 376/23-CSC (SEFAZ-AM),
pode ser impugnada como oferta \textbf{inexequível} nos termos do \art{59}, II,
da Lei n.\textsuperscript{o}~14.133/2021, caso a empresa não demonstre seu
modelo de remuneração indireta. O TR deve exigir que o licitante comprove,
na fase de habilitação econômica, a \textbf{viabilidade financeira da taxa
ofertada}, indicando a receita acessória que sustenta a operação (float
financeiro, comissão das armadoras, etc.).
\end{warnbox}

A sustentabilidade de taxas simbólicas repousa em três fontes de receita
indireta que o TR deve admitir explicitamente e a Administração deve monitorar:

\begin{enumerate}
  \item \textbf{Float financeiro}: a prefeitura paga a fatura com 15 a 30 dias
    de atraso em relação ao embarque; nesse intervalo, a agência retém os
    recursos e obtém rendimento em aplicações financeiras;
  \item \textbf{Comissão das armadoras}: as empresas de navegação remuneram
    canais de venda com 5--10\% do valor da passagem; essa receita não é
    debitada ao Município;
  \item \textbf{Retorno de bilhetes não utilizados (No-Show)}: discutido na
    Seção~\ref{sec:noshow}.
\end{enumerate}

O TR deve estabelecer que a taxa de agenciamento é o \textbf{único valor
licitado} e que qualquer comissão recebida pela agência das armadoras deve ser
declarada ao Gestor sem que isso implique redução do valor da passagem para o
Município.

% ═════════════════════════════════════════════════════════════════════
\section{Dinâmica Hidrológica do Rio Madeira e Impacto nos Custos}
% ═════════════════════════════════════════════════════════════════════

\subsection{Coeficiente de Ajuste Logístico}

O Rio Madeira exibe um regime hidrológico bimodal com amplitude de nível
d'água superior a 12 metros entre a cheia plena e a seca extrema. Esta
amplitude, a maior entre os tributários do Rio Amazonas, produz impactos
diretos sobre a distância efetiva percorrida e o consumo de combustível:

\begin{equation}
  C_{log} = \frac{D_{vazante}}{D_{cheia}} \times \frac{C_{comb}}{C_{base}}
  \label{eq:clog}
\end{equation}

\begin{table}[H]\centering
\caption{Matriz de risco climático e variação de custos --- corredor
  Borba--Manaus}
\small\rowcolors{2}{rowalt}{white}
\begin{tabular}{L{3.2cm} C{2.0cm} C{2.0cm} C{2.2cm} C{2.2cm} C{2.2cm}}
  \rowcolor{navy}
  \color{white}\textbf{Cenário\newline Hidrológico} &
  \color{white}\textbf{Período} &
  \color{white}\textbf{$\Delta$\newline Distância} &
  \color{white}\textbf{$\Delta$\newline Consumo} &
  \color{white}\textbf{Reajuste\newline Provável} &
  \color{white}\textbf{Tarifa Méd.\newline B--M (R\$)} \\
  \toprule
  Cheia plena        & Jan--Mai & $+0\%$  & $+0\%$  & Base     & 320,00 \\
  Vazante moderada   & Jun--Ago & $+8\%$  & $+12\%$ & $+10\%$  & 352,00 \\
  Seca extrema       & Set--Nov & $+15\%$ & $+25\%$ & $+20\%$  & 384,00 \\
  Enchente inicial   & Dez      & $+3\%$  & $+5\%$  & $+3\%$   & 330,00 \\
  \rowcolor{ltblue}
  \multicolumn{4}{l}{\textbf{Média ponderada anual (dias por período)}} &
    $+8{,}5\%$ & \textbf{345,00} \\
  \bottomrule
\end{tabular}
\end{table}

\noindent A média ponderada de \Rs 345,00 por passagem, incorporando a
variabilidade sazonal, fundamenta o valor orçado e substitui a tarifa
estática de \Rs 260,00 do cenário misto, que subestimava os efeitos de
seca extrema observados em 2023 e 2024.

\subsection{Sazonalidade de Demanda e Picos Festivos}

Além da sazonalidade hidrológica, há uma sazonalidade de \textbf{demanda}
que pressiona tanto a disponibilidade de vagas quanto as tarifas de mercado:

\begin{table}[H]\centering
\caption{Sazonalidade de demanda no corredor Borba--Manaus}
\small\rowcolors{2}{rowalt}{white}
\begin{tabular}{C{1.2cm} C{1.8cm} L{4.5cm} C{2.2cm} C{2.0cm}}
  \rowcolor{navy}
  \color{white}\textbf{Mês} &
  \color{white}\textbf{Período\newline Hidrológico} &
  \color{white}\textbf{Evento / Fator} &
  \color{white}\textbf{Fator de\newline Demanda ($\alpha$)} &
  \color{white}\textbf{Correção\newline de Volume} \\
  \toprule
  Jan & Cheia    & Retorno pós-festas               & 1,40 & $+40\%$ \\
  Fev & Cheia    & Carnaval (Manaus)                 & 1,25 & $+25\%$ \\
  Mar & Cheia    & Normal                            & 1,00 & Base \\
  Abr & Cheia    & Semana Santa                      & 1,20 & $+20\%$ \\
  Mai & Cheia    & Normal                            & 1,00 & Base \\
  Jun & Vazante  & Festa de Santo Antônio de Borba
                   (13/06)                           & 1,60 & $+60\%$ \\
  Jul & Vazante  & Férias escolares + seca incipiente& 1,35 & $+35\%$ \\
  Ago & Vazante  & Seca moderada                     & 0,90 & $-10\%$ \\
  Set & Seca     & Seca extrema (risco de cancelamento)& 0,80 & $-20\%$ \\
  Out & Seca     & Início da enchente                & 0,85 & $-15\%$ \\
  Nov & Enchente & Normal                            & 1,00 & Base \\
  Dez & Enchente & Natal / Réveillon                 & 1,50 & $+50\%$ \\
  \bottomrule
\end{tabular}
\end{table}

O TR deverá exigir que a agência \textbf{antecipe reservas} nos meses com
fator $\alpha > 1{,}30$ (janeiro, junho, julho e dezembro), mediante plano
de contingência de embarque apresentado 30 dias antes do evento, sob pena
de multa por indisponibilidade de vagas.

A Festa de Santo Antônio de Borba (13 de junho) merece atenção especial: é
o maior evento cultural do município, com fluxo intensificado nos 5 dias
anteriores e 3 dias posteriores ao dia principal. O pico de $\alpha = 1{,}60$
é o mais crítico de todo o calendário e coincide com o início da vazante,
quando o número de embarcações operando regularmente começa a diminuir.

% ═════════════════════════════════════════════════════════════════════
\section{Modelo de Gestão do No-Show e Créditos de Bilhetes}
\label{sec:noshow}
% ═════════════════════════════════════════════════════════════════════

O \textit{No-Show} --- situação em que o passageiro não comparece ao embarque
--- representa um dos maiores riscos de desperdício de recursos públicos em
contratos de agenciamento. A modelagem matemática do fenômeno é essencial para
dimensionar a perda potencial e definir as cláusulas contratuais adequadas.

\subsection{Modelo Probabilístico de No-Show}

Seja $p_{ns}$ a probabilidade de no-show por bilhete emitido. Com base em
contratos similares (SEMA-AM, SSP-AM, PE 376/23-CSC), estima-se
$p_{ns} \in [0{,}08;\,0{,}15]$ para o transporte fluvial institucional,
onde 8\% é o cenário otimista (com política de confirmação obrigatória 24h
antes) e 15\% o cenário pessimista (sem controle).

A perda financeira esperada anual por no-show não gerenciado é:

\begin{equation}
  L_{ns} = N_{bilhetes} \times p_{ns} \times \bar{P} \times (1 - r_{rec})
  \label{eq:noshow}
\end{equation}

onde $r_{rec}$ é a taxa de recuperação (remarcação ou crédito) que a agência
consegue negociar com o armador. Cenários:

\begin{table}[H]\centering
\caption{Impacto financeiro do No-Show --- cenários anuais}
\small\rowcolors{2}{rowalt}{white}
\begin{tabular}{L{3.8cm} C{1.5cm} C{1.5cm} C{2.5cm} C{2.5cm} C{2.0cm}}
  \rowcolor{navy}
  \color{white}\textbf{Cenário} &
  \color{white}\textbf{$p_{ns}$} &
  \color{white}\textbf{$r_{rec}$} &
  \color{white}\textbf{Bilhetes\newline Perdidos} &
  \color{white}\textbf{Perda Anual\newline (R\$)} &
  \color{white}\textbf{\% do\newline Contrato} \\
  \toprule
  Sem controle algum      & 15\% & 0\%  & 1.628 & 561.660 & 20,1\% \\
  Gestão mínima da agência& 12\% & 40\% & 782   & 269.790 &  9,6\% \\
  Gestão ativa (meta TR)  &  8\% & 75\% & 218   &  75.210 &  2,7\% \\
  \rowcolor{ltblue}
  \multicolumn{3}{l}{\textbf{Meta contratual obrigatória}} &
    \multicolumn{3}{c}{\textbf{$p_{ns} \leq 10\%$ e $r_{rec} \geq 70\%$}} \\
  \bottomrule
\end{tabular}
\end{table}

O TR deve estabelecer como \textbf{KPI mandatório}: taxa de no-show
$\leq 10\%$ e taxa de recuperação $\geq 70\%$ (remarcação ou crédito para
uso em 60 dias). O descumprimento desses índices por 2 meses consecutivos
configurará inexecução parcial e embasará a aplicação de multa.

\subsection{Fluxo de Gestão do No-Show}

A cláusula contratual deve estabelecer o seguinte fluxo:

\begin{enumerate}
  \item O passageiro comunica o no-show à agência com antecedência mínima
    de \textbf{4 horas} antes do horário de embarque;
  \item A agência tenta remarcação em embarcação do mesmo armador, sem ônus
    para o Município;
  \item Impossível a remarcação, gera-se \textbf{crédito} vinculado ao centro
    de custo da secretaria solicitante, válido por \textbf{60 dias corridos};
  \item Crédito não utilizado em 60 dias é convertido em
    \textbf{glosa na fatura} do mês subsequente;
  \item A agência apresenta relatório mensal de no-shows, remarcações e
    créditos, que o Fiscal Técnico concilia com as Ordens de Transporte emitidas.
\end{enumerate}

% ═════════════════════════════════════════════════════════════════════
\section{Análise de Demanda Setorial e Estrutura de Consumo}
% ═════════════════════════════════════════════════════════════════════

\subsection{Escopo do Contrato: Municipal ou Setorial?}

\begin{alertbox}[Decisão Estratégica Prévia ao ETP]
\small
O processo \procNum, instaurado pela \textbf{Secretaria Municipal de Saúde},
pode ser estruturado como:
\begin{itemize}[topsep=2pt,itemsep=1pt]
  \item \textbf{Contrato setorial} (SMS como único contratante):
    valor menor, TR mais simples, menor poder de barganha;
  \item \textbf{ARP municipal} (SMS como órgão gerenciador, demais secretarias
    como participantes): poder de compra consolidado, maior competitividade,
    melhor gestão do float. \textbf{Recomendado.}
\end{itemize}
A opção pela ARP municipal exige que as demais secretarias formalizem
adesão ao processo antes da publicação do edital, mediante DFD setorial
encaminhado à SMS (\art{82}, \S\,3.\textsuperscript{o},
Lei n.\textsuperscript{o}~14.133/2021).
\end{alertbox}

\subsection{Estimativa de Demanda por Secretaria}

A demanda por transporte fluvial institucional em Borba apresenta perfil
\textbf{predominantemente probabilístico}, não determinístico. Os dados
abaixo constituem projeções estatísticas baseadas em:
(i) séries históricas de empenhos da prefeitura (2022--2025);
(ii) média de Tratamentos Fora do Domicílio (TFD) homologados pelo SUS;
(iii) calendário escolar e administrativo.

\begin{table}[H]\centering
\caption{Estimativa de demanda anual por secretaria --- Processo \procNum}
\small\rowcolors{2}{rowalt}{white}
\begin{tabular}{L{3.2cm} L{3.2cm} C{2.0cm} C{2.0cm} C{2.8cm}}
  \rowcolor{navy}
  \color{white}\textbf{Secretaria} &
  \color{white}\textbf{Perfil\newline de Usuário} &
  \color{white}\textbf{Pass./Mês\newline (estimado)} &
  \color{white}\textbf{Pass./Ano} &
  \color{white}\textbf{Custo Anual\newline Estimado (R\$)} \\
  \toprule
  Saúde (SMS)    & Pacientes + acomp. TFD     & 450 & 5.400 & 1.863.000 \\
  Educação (SEMED) & Professores + técnicos   & 200 & 2.400 &   828.000 \\
  Assist. Social & Famílias vulneráveis       & 100 & 1.200 &   414.000 \\
  Administração  & Servidores / diligências   & 150 & 1.800 &   621.000 \\
  \rowcolor{ltblue}
  \textbf{Total} & & \textbf{900} & \textbf{10.800} &
    \textbf{3.726.000} \\
  \bottomrule
\end{tabular}
\end{table}

\begin{alertbox}[Análise do Valor Projetado \textit{vs.} Orçamento]
\small
A projeção ``bruta'' totaliza \Rs 3.726.000, superior ao orçamento de
\Rs 2.800.000. Isso é \textbf{esperado e correto} em uma ARP: o valor
registrado é o \textit{teto de empenho}, não um compromisso de consumo
total. O dimensionamento deve acomodar que: (i) nem todas as passagens
estimadas se concretizarão (fator de realização histórico: 70--80\%);
(ii) a negociação de descontos coletivos pela agência reduzirá o custo
unitário em relação à tarifa ARSEP de balcão.
\end{alertbox}

Aplicando o fator de realização histórico de 75\%:

\begin{align}
  \text{Demanda realizada estimada} &= 10.800 \times 0{,}75 = 8.100
  \text{ passagens/ano} \notag\\
  \text{Custo médio ajustado} &= \Rs 345{,}00 \times 0{,}95
  \text{ (desconto coletivo de 5\%)} = \Rs 327{,}75 \notag\\
  \text{Custo total projetado} &= 8.100 \times 327{,}75 =
  \Rs 2.654.775 \notag
\end{align}

Este valor de \Rs 2.654.775 é compatível com o orçamento de \Rs 2.800.000,
deixando uma reserva de \Rs 145.225 (5,2\%) para o frete de volumes e para
eventuais picos de demanda não previstos.

\subsection{Tratamento Fora do Domicílio (TFD) --- Especificidades}

O TFD é obrigação constitucional de saúde (\art{196} da CF/88) e responde
por aproximadamente 56\% da demanda estimada da SMS. Suas características
operacionais impõem exigências específicas ao agenciamento:

\begin{enumerate}
  \item \textbf{Flexibilidade de data de retorno}: a alta médica ou o
    reagendamento de exame pode alterar a data de volta com menos de 24h de
    antecedência; o contrato deve prever \textbf{remarcação gratuita} em até
    2 ocorrências por paciente por mês;

  \item \textbf{Acompanhante legal}: toda remoção de paciente gera
    \textbf{2 bilhetes} (paciente + acompanhante), nos termos do
    \art{16} da Lei n.\textsuperscript{o}~10.741/2003 (Estatuto do Idoso)
    e do \art{12} da Lei n.\textsuperscript{o}~8.069/1990 (ECA) para
    crianças; o TR deve prever este multiplicador;

  \item \textbf{Urgências e emergências}: o TR deve distinguir o
    \textbf{agenciamento de urgência} (embarque em até 4h da solicitação)
    do \textbf{agenciamento padrão} (embarque programado com 48h de
    antecedência), com SLAs diferentes e possível sobrepreço de taxa de
    urgência limitado a \Rs [X] por bilhete, a ser definido na pesquisa
    de mercado;

  \item \textbf{Pacientes acamados ou cadeirantes}: embarcações de linha
    regular possuem acessibilidade limitada; o TR deve exigir da agência
    a indicação de embarcações adaptadas quando disponíveis, e a
    justificativa por escrito quando não houver.
\end{enumerate}

% ═════════════════════════════════════════════════════════════════════
\section{Estrutura Jurídico-Administrativa sob a Lei 14.133/2021}
% ═════════════════════════════════════════════════════════════════════

\subsection{Modalidade, Instrumento e Código de Classificação}

\begin{table}[H]\centering\small
\caption{Parâmetros jurídico-administrativos do Processo \procNum}
\rowcolors{2}{rowalt}{white}
\begin{tabular}{>{\bfseries\color{navy}}p{4.5cm} L{9.5cm}}
\midrule
Modalidade:                & Pregão Eletrônico (\art{28}, I,
                             Lei n.\textsuperscript{o}~14.133/2021) \\
\rowcolor{rowalt}
Instrumento:               & Ata de Registro de Preços ---
                             ARP (\art{82}) \\
Critério de julgamento:    & Menor taxa de agenciamento por bilhete
                             emitido (\art{33}, I) \\
\rowcolor{rowalt}
Natureza da despesa:       & \textbf{33.90.39} --- Outros Serviços de
                             Terceiros --- Pessoa Jurídica \\
CATSER principal:          & \textbf{3719} --- Agenciamento de Passagens \\
\rowcolor{rowalt}
CATSER secundário:         & \textbf{17892} --- Transporte de Passageiros
                             (frete volumes) \\
Unidade de medida:         & Bilhete emitido (passagem A/V) \\
\rowcolor{rowalt}
Vigência da ARP:           & 12 meses prorrogáveis (\art{84},
                             \S\,1.\textsuperscript{o}) \\
Publicidade obrigatória:   & PNCP + Diário Oficial dos Municípios
                             do Amazonas (\art{174}) \\
\rowcolor{rowalt}
Controle externo:          & TCE-AM + Portal da Transparência Municipal \\
\bottomrule
\end{tabular}
\end{table}

\subsection{Justificativa do SRP para Serviço de Agenciamento}

O Sistema de Registro de Preços é especialmente adequado para este objeto
por quatro razões convergentes:

\begin{enumerate}
  \item \textbf{Demanda imprevisível} (\art{82}, I): o número exato de
    passagens é probabilístico, não determinístico --- o SRP permite empenhos
    parciais conforme a necessidade efetiva, sem o risco de sub ou
    superempenhamento;

  \item \textbf{Entregas parceladas} (\art{82}, II): as passagens são emitidas
    diariamente, em volumes variáveis, ao longo de 12 meses;

  \item \textbf{Conveniência do registro} (\art{82}, IV): múltiplas secretarias
    podem aderir ao mesmo instrumento, eliminando a proliferação de processos
    licitatórios paralelos para o mesmo serviço;

  \item \textbf{Impossibilidade de definir quantidades exatas}: a própria
    natureza do TFD --- dependente de diagnósticos, alta hospitalar e
    agendamentos do SUS --- impede a fixação de quantitativos precisos
    antes da licitação.
\end{enumerate}

% ═════════════════════════════════════════════════════════════════════
\section{Requisitos Técnicos e Regulatórios das Embarcações}
% ═════════════════════════════════════════════════════════════════════

\begin{corrbox}[Correção 1 já aplicada --- DPVAT $\to$ DPEM]
\small
Todos os requisitos de seguro obrigatório referem-se ao
\textbf{DPEM --- Dano Pessoal Causado por Embarcações ou por sua Carga},
regulado pela Resolução CNSP n.\textsuperscript{o}~373/2018 e administrado
pela SUSEP. O DPVAT foi extinto em 2020 (Lei n.\textsuperscript{o}~14.071/2020)
e não tem aplicação para embarcações.
\end{corrbox}

\begin{table}[H]\centering
\caption{Requisitos técnicos mínimos para embarcações cadastradas pela agência}
\small\rowcolors{2}{rowalt}{white}
\begin{tabular}{L{4.5cm} L{4.0cm} L{4.0cm} L{1.5cm}}
  \rowcolor{navy}
  \color{white}\textbf{Requisito} &
  \color{white}\textbf{Lancha Rápida (A-Jato)} &
  \color{white}\textbf{Barco Motor (Recreio)} &
  \color{white}\textbf{Base Legal} \\
  \toprule
  Registro ANTAQ   & Obrigatório & Obrigatório & Res. ANTAQ 22/2013 \\
  Autori. ARSEP-AM & Obrigatório & Obrigatório & Dec. Est. AM 34.536/14 \\
  DPEM (SUSEP)     & Obrigatório & Obrigatório & CNSP Res. 373/2018 \\
  Salva-vidas      & 110\% da lotação & 110\% da lotação & NORMAM-02/DPC \\
  Certificado de\newline Segurança &
    Marinha do Brasil (DPC) & Marinha do Brasil (DPC) & NORMAM-01/DPC \\
  Arq. de Bordagem & Obrigatório & Obrigatório & CONTRAN/ANTAQ \\
  Ar condicionado  & Padrão & Opcional (suítes) & TR \\
  Velocidade de\newline cruzeiro &
    25--35 nós & 8--12 nós & Projeto naval \\
  \bottomrule
\end{tabular}
\end{table}

A agência deverá manter e apresentar mensalmente ao Gestor:

\begin{itemize}
  \item Banco de dados atualizado com nome, identificação ANTAQ, registro
    ARSEP e apólice DPEM de cada embarcação utilizada;
  \item Relatório de embarcações com documentação vencida, com prazo de 5 dias
    para substituição;
  \item Lista de embarcações adaptadas para mobilidade reduzida, quando
    disponíveis no corredor.
\end{itemize}

% ═════════════════════════════════════════════════════════════════════
\section{Fluxo Operacional: Ordem de Transporte}
% ═════════════════════════════════════════════════════════════════════

O instrumento operacional que substitui a ``compra no guichê'' é a
\textbf{Ordem de Transporte (OT)}, documento emitido pelo Gestor ou por
servidor autorizado de cada secretaria, que aciona a agência para emissão
do(s) bilhete(s). O fluxo deve ser detalhado no TR:

\begin{table}[H]\centering\small
\caption{Fluxo operacional da Ordem de Transporte}
\rowcolors{2}{rowalt}{white}
\begin{tabular}{C{0.6cm} L{4.0cm} L{4.5cm} C{2.0cm} L{2.5cm}}
  \rowcolor{navy}
  \color{white}\textbf{Etapa} &
  \color{white}\textbf{Ação} &
  \color{white}\textbf{Responsável} &
  \color{white}\textbf{Prazo} &
  \color{white}\textbf{Documento} \\
  \toprule
  1 & Solicitação de viagem com dados do passageiro &
    Secretaria solicitante & $D-48h$ (padrão) ou $D-4h$ (urgência) &
    OT via sistema ou e-mail \\
  2 & Verificação de disponibilidade e tarifa &
    Agência & 1h após OT & Confirmação de reserva \\
  3 & Emissão do bilhete eletrônico (QR Code) &
    Agência & 2h após confirmação & Bilhete digital (e-mail/app) \\
  4 & Confirmação de embarque &
    Agência / armador & No embarque & Manifesto digital \\
  5 & Relatório mensal consolidado &
    Agência & Até 5\textdegree\ dia útil & Fatura + manifesto \\
  6 & Conciliação e atestado &
    Fiscal Técnico & Até 10.\textdegree\ dia útil &
    TRD (equivalente) \\
  7 & Pagamento &
    Setor financeiro & Até 30 dias após fatura & OB/TED/PIX \\
  \bottomrule
\end{tabular}
\end{table}

\noindent A OT deve conter, no mínimo: nome completo e CPF do passageiro;
secretaria de origem e centro de custo; trecho (ida/volta ou só ida); data e
horário preferencial de embarque; modal preferencial (justificar se jato);
natureza da viagem (TFD, serviço, diligência, outro).

% ═════════════════════════════════════════════════════════════════════
\section{Requisitos Tecnológicos da Agência}
% ═════════════════════════════════════════════════════════════════════

A contratada deve disponibilizar plataforma tecnológica com as seguintes
funcionalidades, como condição de habilitação técnica:

\begin{enumerate}
  \item \textbf{Portal de self-booking} ou canal de atendimento 24/7 via
    aplicativo de mensagens (WhatsApp Business API ou equivalente), dado
    que embarcações de linha regular em Borba frequentemente partem ou
    chegam em horários noturnos;

  \item \textbf{Bilhete eletrônico} com QR Code, nome do passageiro e
    número da OT gravado no código, permitindo auditoria pelo TCE-AM;

  \item \textbf{Dashboard de gestão} acessível ao Fiscal Técnico e ao Gestor,
    com visão em tempo real de: bilhetes emitidos por secretaria, créditos
    abertos por no-show, saldo de ARP por lote e KPIs do contrato;

  \item \textbf{Integração contábil}: exportação de relatório mensal em
    formato compatível com o sistema de contabilidade do Município
    (planilha CSV/XLSX), com campos de centro de custo, número da OT,
    CPF do passageiro e valor por bilhete, para conciliação com o SICONFI;

  \item \textbf{Backup de contingência}: canal alternativo de atendimento
    (telefone fixo) para operação em caso de falha do sistema digital,
    considerando a conectividade intermitente de Borba.
\end{enumerate}

% ═════════════════════════════════════════════════════════════════════
\section{Fórmula de Reajuste e Reequilíbrio Econômico-Financeiro}
% ═════════════════════════════════════════════════════════════════════

\subsection{Reajuste Paramétrico Anual}

O contrato de agenciamento está exposto a dois vetores inflacionários
distintos: o custo operacional da agência (capturado pelo IPCA) e o custo
do combustível das embarcações (que, no interior do Amazonas, supera a
inflação geral pelo custo do frete de abastecimento por balsa).

A fórmula paramétrica proposta pondera esses dois vetores:

\begin{equation}
  \hat{R} = 0{,}60\left(\frac{D_1}{D_0}\right)
           + 0{,}40\left(\frac{IPCA_1}{IPCA_0}\right)
  \label{eq:reajuste}
\end{equation}

onde $D_1/D_0$ é a variação do preço médio do óleo diesel marítimo pesquisado
pela ANP no Estado do Amazonas no período de referência, e $IPCA_1/IPCA_0$
a variação do índice oficial de inflação no mesmo período. O peso de 60\%
para o diesel reflete sua participação de 45--60\% na estrutura de custos
da navegação amazônica.

\begin{table}[H]\centering
\caption{Simulação de reajuste em três cenários --- 12 meses}
\small\rowcolors{2}{rowalt}{white}
\begin{tabular}{L{3.5cm} C{2.0cm} C{2.0cm} C{2.5cm} C{3.0cm}}
  \rowcolor{navy}
  \color{white}\textbf{Cenário} &
  \color{white}\textbf{$\Delta$ Diesel\newline (\%)} &
  \color{white}\textbf{IPCA\newline (\%)} &
  \color{white}\textbf{$\hat{R}$\newline (\%)} &
  \color{white}\textbf{Impacto no\newline Contrato (R\$)} \\
  \toprule
  Inflação baixa       &  $+5\%$  & $+4\%$  &  4,6\% &  128.800 \\
  Cenário base         & $+10\%$  & $+5\%$  &  8,0\% &  224.000 \\
  Choque de combustível& $+25\%$  & $+7\%$  & 17,8\% &  498.400 \\
  \bottomrule
\end{tabular}
\end{table}

\noindent O reajuste somente incide sobre a taxa de agenciamento por bilhete
(parcela fixa), \textbf{não sobre o valor da passagem}, que é reajustado
autonomamente pela ARSEP-AM via revisão tarifária.

\subsection{Reequilíbrio por Alteração Tarifária da ARSEP-AM}

Se a ARSEP-AM publicar nova tabela de tarifas durante a vigência da ARP, o
Município não poderá ser onerado além do registrado. A cláusula de
reequilíbrio deve prever:

\begin{enumerate}
  \item \textbf{Aumento tarifário}: a agência absorve a diferença por 30 dias
    e, ao fim desse período, apresenta pedido de revisão fundamentado com
    portaria da ARSEP-AM; o Gestor tem 15 dias para autorizar o novo preço
    de passagem a ser repassado;
  \item \textbf{Redução tarifária}: o benefício é imediatamente repassado ao
    Município; vedada qualquer retenção pela agência.
\end{enumerate}

% ═════════════════════════════════════════════════════════════════════
\section{Matriz de Riscos Expandida}
% ═════════════════════════════════════════════════════════════════════

\begin{table}[H]\centering
\caption{Matriz de riscos --- contratação de agenciamento fluvial}
\small\rowcolors{2}{rowalt}{white}
\begin{tabular}{C{0.6cm} L{3.8cm} C{1.3cm} C{1.3cm} C{1.3cm} L{4.5cm}}
  \rowcolor{navy}
  \color{white}\textbf{R\#} &
  \color{white}\textbf{Risco} &
  \color{white}\textbf{Prob.\newline (1--5)} &
  \color{white}\textbf{Impacto\newline (1--5)} &
  \color{white}\textbf{Score} &
  \color{white}\textbf{Mitigação} \\
  \toprule
  R1 & Seca extrema reduz oferta de embarcações & 4 & 5 & 20 &
    Cadastro mínimo de 5 armadores por modal; cláusula de contingência \\
  R2 & Agência em dificuldade financeira
    (insolvência / interrupção) & 2 & 5 & 10 &
    Exigir garantia contratual de 5\% (\art{96}, Lei 14.133/2021);
    monitorar certidões mensalmente \\
  R3 & Pico de demanda em Santo Antônio /
    Natal sem vagas disponíveis & 4 & 4 & 16 &
    Plano de contingência 30 dias antes; reserva antecipada obrigatória \\
  R4 & No-show elevado ($> 15\%$) & 3 & 3 & 9 &
    KPI contratual + glosa automática de créditos não utilizados \\
  R5 & Reajuste tarifário ARSEP-AM não previsto & 3 & 4 & 12 &
    Cláusula de reequilíbrio automático (Seção~8.2) \\
  R6 & Embarcação sem DPEM vigente utilizada
    para passageiro público & 2 & 5 & 10 &
    Agência reporta apólices mensalmente; Fiscal verifica in loco \\
  R7 & Choque do preço do diesel (+30\%) & 3 & 4 & 12 &
    Fórmula paramétrica $\hat{R}$ com gatilho automático \\
  R8 & Falha de sistema / bilhete não gerado
    em urgência & 3 & 4 & 12 &
    SLA de contingência: bilhete em papel assinado pelo gerente,
    digitalizado em 24h \\
  R9 & Conluio entre armadoras para elevar tarifa
    acima da ARSEP & 2 & 4 & 8 &
    Agência cadastra $\geq 5$ armadores; monitoramento mensal de tarifas \\
  R10 & Cancelamento de viagem por acidente
    fluvial & 2 & 5 & 10 &
    Seguro DPEM cobre passageiros; agência tem 4h para reacomodar \\
  \bottomrule
\end{tabular}
\end{table}

% ═════════════════════════════════════════════════════════════════════
\section{Indicadores de Desempenho (KPIs) do Contrato}
% ═════════════════════════════════════════════════════════════════════

\begin{table}[H]\centering
\caption{KPIs propostos para o contrato de agenciamento ---
  Processo \procNum}
\small\rowcolors{2}{rowalt}{white}
\begin{tabular}{L{4.5cm} L{4.0cm} C{2.0cm} C{2.0cm}}
  \rowcolor{navy}
  \color{white}\textbf{Indicador} &
  \color{white}\textbf{Fórmula} &
  \color{white}\textbf{Meta} &
  \color{white}\textbf{Freq.} \\
  \toprule
  Taxa de no-show &
    Bilhetes sem embarque / Total emitido & $\leq 10\%$ & Mensal \\
  Taxa de recuperação (remarcação) &
    Remarcações / No-shows & $\geq 70\%$ & Mensal \\
  Tempo de emissão --- padrão &
    Tempo entre OT e bilhete digital (h) & $\leq 2h$ & Mensal \\
  Tempo de emissão --- urgência &
    Tempo entre OT urgência e bilhete (h) & $\leq 1h$ & Mensal \\
  Conformidade documental das embarc. &
    Embarcações c/ DPEM e ANTAQ válidos /
    Total embarcações & $= 100\%$ & Mensal \\
  Disponibilidade em picos ($\alpha > 1{,}3$) &
    OT atendidas em picos / OT solicitadas & $\geq 95\%$ & Por evento \\
  Tempestividade do relatório mensal &
    Entregue até 5.\textdegree\ d.u. do mês & $= 100\%$ & Mensal \\
  Satisfação do usuário (escala 1--5) &
    Pesquisa de satisfação semestral & $\geq 4{,}0$ & Semestral \\
  \bottomrule
\end{tabular}
\end{table}

% ═════════════════════════════════════════════════════════════════════
\section{Conformidade Regulatória e Papel da ARSEP-AM}
% ═════════════════════════════════════════════════════════════════════

A Agência Reguladora de Serviços Públicos Delegados e Contratados do Estado
do Amazonas (ARSEP-AM) exerce competência sobre o transporte aquaviário de
passageiros intermunicipal, definindo as tarifas máximas para o corredor
Borba--Manaus e as condições operacionais mínimas das embarcações autorizadas.

A agência contratada pelo Município deverá:

\begin{enumerate}
  \item Contratar \textbf{exclusivamente} embarcações detentoras de
    Autorização de Transporte de Passageiros emitida pela ARSEP-AM;

  \item Respeitar as tarifas máximas da tabela ARSEP-AM vigente ao momento
    de emissão de cada bilhete; vedada a prática de tarifas acima do teto
    regulado;

  \item Comunicar ao Gestor, no prazo de 48h, qualquer alteração publicada
    pela ARSEP-AM na tabela tarifária do corredor Borba--Manaus;

  \item Não contratar embarcações em situação de operação clandestina
    (sem autorização ARSEP/ANTAQ), ainda que ofereçam tarifas menores;
    essa prática expõe o Município a responsabilidade solidária por danos
    a passageiros.
\end{enumerate}

% ═════════════════════════════════════════════════════════════════════
\section{Governança e Vantagens da Centralização}
% ═════════════════════════════════════════════════════════════════════

A contratação de empresa de agenciamento, em vez da compra direta pulverizada,
produz as seguintes vantagens mensuráveis para a governança municipal:

\begin{table}[H]\centering
\caption{Comparativo de modelos: agenciamento \textit{vs.} compra direta}
\small\rowcolors{2}{rowalt}{white}
\begin{tabular}{L{4.5cm} L{4.5cm} L{4.5cm}}
  \rowcolor{navy}
  \color{white}\textbf{Dimensão} &
  \color{white}\textbf{Agenciamento (recomendado)} &
  \color{white}\textbf{Compra direta} \\
  \toprule
  Auditoria pelo TCE-AM &
    Fatura única com manifesto e QR Code; rastreável por CPF &
    Reembolsos com notas manuais; difícil conciliação \\
  Controle por secretaria &
    Dashboard por centro de custo em tempo real &
    Planilhas manuais; reconciliação mensal \\
  Negociação de tarifas &
    Desconto coletivo de 5--15\% sobre tabela ARSEP &
    Tarifa de balcão; sem poder de barganha \\
  Gestão de no-show &
    Sistema de crédito e remarcação automática &
    Passagem perdida sem recuperação \\
  Processos de dispensa &
    Eliminados; um único contrato anual &
    Dezenas de dispensas de pequena monta \\
  Risco de superfaturamento &
    Controlado pela fatura da agência + manifesto &
    Alto risco (reembolso sem comprovação de embarque) \\
  Carga administrativa &
    Baixa (Gestor valida fatura mensal) &
    Alta (Gestor processa cada reembolso) \\
  \bottomrule
\end{tabular}
\end{table}

% ═════════════════════════════════════════════════════════════════════
\section{Síntese e Conclusão: Viabilidade do Orçamento de
  \Rs \valorTotal}
% ═════════════════════════════════════════════════════════════════════

O presente estudo valida, com a correção e o detalhamento necessários, o
montante de \textbf{\Rs \valorTotal} como adequado e suficiente para um
ciclo anual de agenciamento de transporte fluvial no corredor Borba--Manaus,
com cobertura das quatro principais secretarias do Município.

A validade repousa em quatro premissas verificáveis:

\begin{enumerate}
  \item Demanda realizada de \textbf{8.100 passagens/ano} (fator de
    realização 75\% sobre a projeção bruta), compatível com o porte
    demográfico de Borba (34.869 hab.) e com a estrutura do TFD-SUS;

  \item Custo médio ponderado de \textbf{\Rs 327,75 por passagem}
    (tarifa ARSEP com desconto coletivo de 5\%), coerente com os praticados
    no corredor em 2024--2025;

  \item Reserva de \textbf{5\% do orçamento} (\Rs 140.000) para frete de
    volumes e eventos de pico não previstos;

  \item Fórmula paramétrica de reajuste com gatilho no diesel que protege
    a Administração de revisões orçamentárias no ano seguinte.
\end{enumerate}

O risco residual mais relevante é o \textbf{R1} (seca extrema), com score 20,
que deverá ser integralmente absorvido pela cláusula de contingência e pelo
plano de embarque em picos. Os demais riscos têm score $\leq 16$ e são
mitigáveis pelas cláusulas do TR e pelos KPIs contratuais.

\vspace{1.0cm}
\noindent\hfill Borba/AM, \exercicio.

\vspace{1.2cm}
\begin{minipage}{0.47\linewidth}
\begin{center}
  \rule{7.5cm}{0.5pt}\\[3pt]
  \textbf{[Nome e Cargo]}\\
  \small Elaborador(a) da ATP\\
  \small Secretaria Municipal de Saúde de Borba/AM
\end{center}
\end{minipage}
\hfill
\begin{minipage}{0.47\linewidth}
\begin{center}
  \rule{7.5cm}{0.5pt}\\[3pt]
  \textbf{[Nome e Cargo]}\\
  \small Secretário(a) Municipal de Saúde\\
  \small Aprovação
\end{center}
\end{minipage}

\vspace{1.2cm}
\noindent\textcolor{navy!30}{\rule{\linewidth}{0.5pt}}

\noindent{\footnotesize\itshape
\textbf{Base normativa e referências técnicas:}
Lei n.\textsuperscript{o}~14.133/2021 \textbullet{}
Decreto n.\textsuperscript{o}~11.462/2023 \textbullet{}
IN SEGES/MGI n.\textsuperscript{o}~58/2022 \textbullet{}
IN SEGES/MGI n.\textsuperscript{o}~65/2021 \textbullet{}
CNSP Resolução n.\textsuperscript{o}~373/2018 (DPEM) \textbullet{}
Lei n.\textsuperscript{o}~14.071/2020 (extinção do DPVAT) \textbullet{}
NORMAM-01/DPC e NORMAM-02/DPC (Marinha do Brasil) \textbullet{}
Res. ANTAQ n.\textsuperscript{o}~22/2013 \textbullet{}
Dec. Est. AM n.\textsuperscript{o}~34.536/2014 (ARSEP-AM) \textbullet{}
Lei n.\textsuperscript{o}~10.741/2003 (Estatuto do Idoso) \textbullet{}
Lei n.\textsuperscript{o}~8.069/1990 (ECA) \textbullet{}
CF/88, arts.\,37 e 196 \textbullet{}
PE 376/23-CSC (SEFAZ-AM) --- referência de mercado.}

\end{document}
